% GENERATED FILE. Edit canonical lessons, then run npm run generate:books.
% canonical-source-hash: fnv1a64:f6cb50429b2c9253
% canonical-lessons: ML-C537-talaccoru, ML-C537-kutal, ML-C537-kannunir, ML-C537-kurkkamvalikkuka, ML-C537-viyarkkuka

\chapter{Inside the Body}
\label{ch:ml-inside-the-body}
\hlchaptermodality{537}

\begin{chapteropening}
\textbf{By the end of this chapter:} \emph{I can say the brain, the intestines, tears, to snore and to perspire.}

Complete the last lesson of chapter 537: I can say the brain, the intestines, tears, to snore and to perspire.
\end{chapteropening}

\section[\texorpdfstring{\ml{തലച്ചോറ്} (talaccōṟŭ)}{തലച്ചോറ് (talaccōṟŭ)}]{\ml{തലച്ചോറ്} (talaccōṟŭ) — the brain}

\label{lesson:ML-C537-talaccoru}



\begin{quote}
Before the new one: say the Malayalam for the lungs, then the Malayalam for to take a rest.
\end{quote}

\subsection*{You'll want to know: \ml{തലച്ചോറ്}}
\textbf{\ml{തലച്ചോറ്}} — \emph{talaccōṟŭ} — \textquotedblleft{}the brain\textquotedblright{}.

\begin{quote}
\ml{തലച്ചോറിന്} \ml{നല്ല} \ml{ഉറക്കം} \ml{ആവശ്യമാണ്}.
\end{quote}

\begin{quote}
\emph{talaccōṟinŭ nalla uṟakkaṁ āvaśyamāṇŭ.}
\end{quote}

\begin{quote}
The brain needs good sleep.
\end{quote}

\ml{തല} head + \ml{ചോറ്} cooked rice: the soft matter inside the head.

\begin{cousinweb}
\textbf{\ml{തലച്ചോറ്}} (\emph{talaccōṟŭ}) is native Dravidian, a compound of \ml{തല}, \textquotedblleft{}head\textquotedblright{} (DEDR 3103), and \ml{ചോറ്}, which DEDR gives as boiled rice and also brain or marrow (DEDR 2897).
\end{cousinweb}

\begin{grammarlens}[title={What you've built}]
One more word for describing what you lived and felt.
\end{grammarlens}

\subsection*{Your turn}
\begin{itemize}
\raggedright
  \item \emph{Say it:} \emph{talaccōṟŭ}
  \item \emph{Say it:} \emph{talaccōṟŭ}, once more
  \item \emph{Say it:} \emph{talaccōṟinŭ nalla uṟakkaṁ āvaśyamāṇŭ.} — The brain needs good sleep.
  \item \emph{Recall:} say the Malayalam for the lungs, then the Malayalam for to take a rest, then say \emph{talaccōṟŭ} again
\end{itemize}

\begin{tcolorbox}[breakable,colback=teal!4,colframe=teal!35!black,title={Before you move on}]
What does \ml{തലച്ചോറ്} mean? (\textquotedblleft{}The brain\textquotedblright{}.) Say it once more.
\end{tcolorbox}
\section[\texorpdfstring{\ml{കുടൽ} (kuṭal)}{കുടൽ (kuṭal)}]{\ml{കുടൽ} (kuṭal) — the intestines; the gut}

\label{lesson:ML-C537-kutal}



\begin{quote}
Before the new one: say the Malayalam for to take a rest, then the Malayalam for the brain.
\end{quote}

\subsection*{You'll want to know: \ml{കുടൽ}}
\textbf{\ml{കുടൽ}} — \emph{kuṭal} — \textquotedblleft{}the intestines; the gut\textquotedblright{}.

\begin{quote}
\ml{പഴകിയ} \ml{ഭക്ഷണം} \ml{കുടലിന്} \ml{ദോഷം} \ml{ചെയ്യും}.
\end{quote}

\begin{quote}
\emph{paḻakiya bhakṣaṇaṁ kuṭalinŭ dōṣaṁ ceyyuṁ.}
\end{quote}

\begin{quote}
Stale food harms the gut.
\end{quote}

\ml{വൻകുടൽ} is the large intestine, \ml{ചെറുകുടൽ} the small one.

\begin{cousinweb}
\textbf{\ml{കുടൽ}} (\emph{kuṭal}) is native Dravidian: DEDR 1652 gives the Malayalam word as the bowels (DEDR 1652).
\end{cousinweb}

\begin{grammarlens}[title={What you've built}]
One more word for describing what you lived and felt.
\end{grammarlens}

\subsection*{Your turn}
\begin{itemize}
\raggedright
  \item \emph{Say it:} \emph{kuṭal}
  \item \emph{Say it:} \emph{kuṭal}, once more
  \item \emph{Say it:} \emph{paḻakiya bhakṣaṇaṁ kuṭalinŭ dōṣaṁ ceyyuṁ.} — Stale food harms the gut.
  \item \emph{Recall:} say the Malayalam for to take a rest, then the Malayalam for the brain, then say \emph{kuṭal} again
\end{itemize}

\begin{tcolorbox}[breakable,colback=teal!4,colframe=teal!35!black,title={Before you move on}]
What does \ml{കുടൽ} mean? (\textquotedblleft{}The intestines; the gut\textquotedblright{}.) Say it once more.
\end{tcolorbox}
\section[\texorpdfstring{\ml{കണ്ണുനീർ} (kaṇṇunīr)}{കണ്ണുനീർ (kaṇṇunīr)}]{\ml{കണ്ണുനീർ} (kaṇṇunīr) — tears}

\label{lesson:ML-C537-kannunir}



\begin{quote}
Before the new one: say the Malayalam for the brain, then the Malayalam for the intestines.
\end{quote}

\subsection*{You'll want to know: \ml{കണ്ണുനീർ}}
\textbf{\ml{കണ്ണുനീർ}} — \emph{kaṇṇunīr} — \textquotedblleft{}tears\textquotedblright{}.

\begin{quote}
\ml{സിനിമയുടെ} \ml{അവസാനം} \ml{അമ്മയുടെ} \ml{കണ്ണുനീർ} \ml{ഒഴുകി}.
\end{quote}

\begin{quote}
\emph{sinimayuṭe avasānaṁ ammayuṭe kaṇṇunīr oḻuki.}
\end{quote}

\begin{quote}
At the end of the film, Mother's tears flowed.
\end{quote}

\ml{കണ്ണ്} eye + \ml{നീര്} water. \ml{കണ്ണീർ} is a shorter form.

\begin{cousinweb}
\textbf{\ml{കണ്ണുനീർ}} (\emph{kaṇṇunīr}) is native Dravidian, a compound of \ml{കണ്ണ്}, \textquotedblleft{}eye\textquotedblright{} (DEDR 1159), and \ml{നീര്}, \textquotedblleft{}water\textquotedblright{} (DEDR 3690).
\end{cousinweb}

\begin{grammarlens}[title={What you've built}]
One more word for describing what you lived and felt.
\end{grammarlens}

\subsection*{Your turn}
\begin{itemize}
\raggedright
  \item \emph{Say it:} \emph{kaṇṇunīr}
  \item \emph{Say it:} \emph{kaṇṇunīr}, once more
  \item \emph{Say it:} \emph{sinimayuṭe avasānaṁ ammayuṭe kaṇṇunīr oḻuki.} — At the end of the film, Mother's tears flowed.
  \item \emph{Recall:} say the Malayalam for the brain, then the Malayalam for the intestines, then say \emph{kaṇṇunīr} again
\end{itemize}

\begin{tcolorbox}[breakable,colback=teal!4,colframe=teal!35!black,title={Before you move on}]
What does \ml{കണ്ണുനീർ} mean? (\textquotedblleft{}Tears\textquotedblright{}.) Say it once more.
\end{tcolorbox}
\section[\texorpdfstring{\ml{കൂർക്കംവലിക്കുക} (kūrkkaṁvalikkuka)}{കൂർക്കംവലിക്കുക (kūrkkaṁvalikkuka)}]{\ml{കൂർക്കംവലിക്കുക} (kūrkkaṁvalikkuka) — to snore}

\label{lesson:ML-C537-kurkkamvalikkuka}



\begin{quote}
Before the new one: say the Malayalam for the intestines, then the Malayalam for tears.
\end{quote}

\subsection*{You'll want to know: \ml{കൂർക്കംവലിക്കുക}}
\textbf{\ml{കൂർക്കംവലിക്കുക}} — \emph{kūrkkaṁvalikkuka} — \textquotedblleft{}to snore\textquotedblright{}.

\begin{quote}
\ml{അച്ഛൻ} \ml{രാത്രി} \ml{ഉറക്കെ} \ml{കൂർക്കംവലിച്ചു}.
\end{quote}

\begin{quote}
\emph{acchan rātri uṟakke kūrkkaṁvaliccu.}
\end{quote}

\begin{quote}
Father snored loudly in the night.
\end{quote}

Past: \ml{കൂർക്കംവലിച്ചു}, \emph{kūrkkaṁvaliccu}.

\ml{കൂർക്കം} snoring + \ml{വലിക്കുക} to pull, to draw in.

\begin{cousinweb}
\textbf{\ml{കൂർക്കംവലിക്കുക}} (\emph{kūrkkaṁvalikkuka}) is native Dravidian, a compound of \ml{കൂർക്കം}, \textquotedblleft{}snoring\textquotedblright{} (DEDR 1796), and \ml{വലിക്കുക}, \textquotedblleft{}to pull\textquotedblright{} (DEDR 5282).
\end{cousinweb}

\begin{grammarlens}[title={What you've built}]
One more word for describing what you lived and felt.
\end{grammarlens}

\subsection*{Your turn}
\begin{itemize}
\raggedright
  \item \emph{Say it:} \emph{kūrkkaṁvalikkuka}
  \item \emph{Say it:} \emph{kūrkkaṁvalikkuka}, once more
  \item \emph{Say it:} \emph{acchan rātri uṟakke kūrkkaṁvaliccu.} — Father snored loudly in the night.
  \item \emph{Recall:} say the Malayalam for the intestines, then the Malayalam for tears, then say \emph{kūrkkaṁvalikkuka} again
\end{itemize}

\begin{tcolorbox}[breakable,colback=teal!4,colframe=teal!35!black,title={Before you move on}]
What does \ml{കൂർക്കംവലിക്കുക} mean? (\textquotedblleft{}To snore\textquotedblright{}.) Say it once more.
\end{tcolorbox}
\section[\texorpdfstring{\ml{വിയർക്കുക} (viyarkkuka)}{വിയർക്കുക (viyarkkuka)}]{\ml{വിയർക്കുക} (viyarkkuka) — to perspire; to sweat}

\label{lesson:ML-C537-viyarkkuka}



\begin{quote}
Before the new one: say the Malayalam for tears, then the Malayalam for to snore.
\end{quote}

\subsection*{You'll want to know: \ml{വിയർക്കുക}}
\textbf{\ml{വിയർക്കുക}} — \emph{viyarkkuka} — \textquotedblleft{}to perspire; to sweat\textquotedblright{}.

\begin{quote}
\ml{ചൂടിൽ} \ml{ജോലി} \ml{ചെയ്ത്} \ml{തൊഴിലാളികൾ} \ml{വിയർത്തു}.
\end{quote}

\begin{quote}
\emph{cūṭil jōli ceytŭ toḻilāḷikaḷ viyarttu.}
\end{quote}

\begin{quote}
The workers sweated, working in the heat.
\end{quote}

Past: \ml{വിയർത്തു}, \emph{viyarttu}.

The noun \ml{വിയർപ്പ്}, sweat, is already familiar; the past is \ml{വിയർത്തു}.

\begin{cousinweb}
\textbf{\ml{വിയർക്കുക}} (\emph{viyarkkuka}) is native Dravidian, the verb beside \ml{വിയർപ്പ്}, \textquotedblleft{}sweat\textquotedblright{} (DEDR 5486).
\end{cousinweb}

\begin{grammarlens}[title={What you've built}]
One more word for describing what you lived and felt.
\end{grammarlens}

\subsection*{Your turn}
\begin{itemize}
\raggedright
  \item \emph{Say it:} \emph{viyarkkuka}
  \item \emph{Say it:} \emph{viyarkkuka}, once more
  \item \emph{Say it:} \emph{cūṭil jōli ceytŭ toḻilāḷikaḷ viyarttu.} — The workers sweated, working in the heat.
  \item \emph{Recall:} say the Malayalam for tears, then the Malayalam for to snore, then say \emph{viyarkkuka} again
\end{itemize}

\begin{tcolorbox}[breakable,colback=teal!4,colframe=teal!35!black,title={Before you move on}]
What does \ml{വിയർക്കുക} mean? (\textquotedblleft{}To perspire; to sweat\textquotedblright{}.) Say it once more.
\end{tcolorbox}
